\section{Cross-Layer Interface Gaps}
\label{sec:gaps}

The reviewed approaches use different behavioral models, so their guarantees
cannot be combined simply by listing them.  A cross-layer execution argument
must preserve the assumptions that matter when an output expressed under one
model becomes an input to another.  \Cref{tab:coverage} codes each approach
class against seven interface obligations, and four patterns stand out.  No
class is explicit and central on more than two of the seven obligations.  No
class treats membership, ownership, and lifecycle transitions as explicit and
central.  Temporal validity is explicit only in the two distributed-state
classes, which treat physical validity and recovery at most partially.  Recovery
is explicit only where authority transfer is local to one controller or one
runtime.  The obligations group into five recurring gaps, G1--G5.

\begin{table}[!t]
  \centering
  \caption{Qualitative Coverage of Interface Obligations by Approach Class}
  \label{tab:coverage}
  \footnotesize
  \setlength{\tabcolsep}{2.4pt}
  \newcommand{\colhead}[1]{\rotatebox{90}{#1}}
  \begin{tabular}{@{}>{\raggedright\arraybackslash}p{0.5\columnwidth}*{7}{c}@{}}
    \toprule
    & \multicolumn{2}{c}{G1} & G2 & G3 & \multicolumn{2}{c}{G4} & G5 \\
    \cmidrule(lr){2-3}\cmidrule(lr){4-4}\cmidrule(lr){5-5}\cmidrule(lr){6-7}\cmidrule(l){8-8}
    Approach class & \colhead{Semantic} & \colhead{Physical} & \colhead{Temporal}
      & \colhead{Uncertainty} & \colhead{Compatibility} & \colhead{Lifecycle}
      & \colhead{Recovery} \\
    \midrule
    Mission contracts                    & E  & P  & -- & -- & -- & P  & P  \\
    Language-conditioned planning        & E  & P  & -- & -- & -- & -- & P  \\
    Coordinated-mission assurance        & E  & P  & -- & E  & -- & P  & P  \\
    Component contracts and monitors     & P  & P  & P  & -- & P  & -- & P  \\
    Distributed programming and design   & E  & P  & P  & P  & E  & P  & P  \\
    Dynamical contracts                  & -- & E  & -- & -- & -- & -- & -- \\
    Assume--guarantee learning           & P  & P  & -- & E  & -- & -- & -- \\
    Middleware and QoS                   & -- & -- & E  & -- & E  & P  & -- \\
    Communication-aware coordination     & P  & P  & E  & P  & -- & P  & P  \\
    Runtime assurance and safety filters & -- & E  & P  & P  & -- & -- & E  \\
    Mission-level runtime verification   & E  & P  & P  & P  & -- & -- & P  \\
    Cross-layer guardrails and runtimes  & P  & E  & P  & P  & P  & P  & E  \\
    \bottomrule
  \end{tabular}

  \vspace{3pt}
  \parbox{\columnwidth}{\scriptsize E: explicit and central in every source of
  the class; P: partial, conditional, model-internal, or present in only some
  sources; --: outside the stated mechanisms.  Semantic: task identity, intent,
  and hard constraints.  Physical: feasibility or safety under a physical model.
  Temporal: event time, data age, validity interval, and clock semantics.
  Uncertainty: uncertainty representation, provenance, and evidence kind.
  Compatibility: implementation, version, and configuration identity.
  Lifecycle: membership, ownership, and lifecycle transitions.  Recovery:
  defined response to a failed check, including authority transfer.  G1--G5
  refer to the gaps of \cref{sec:gaps}; sources per class are listed in
  \cref{tab:taxonomy}.}
\end{table}


\subsection{G1: Semantic and Physical Validity}

Mission contracts make goals and contexts explicit, and planning and learning
systems produce assignments or actions.  The receiving component still needs the
task identifiers, capabilities, map, state, dynamics, resources, and
communication model behind a proposal.  It must also be able to distinguish a
hard constraint from a planner preference, a model suggestion, and a checked
feasibility claim.  A meaningful plan may be physically infeasible, and a
feasible plan may violate the operator's intent.  Feasibility must therefore be
qualified as nominal, robust, probabilistic, or conditional on a named model and
operating region
\cite{mallozzi2020crome,mallozzi2023repair,mandi2024roco,obata2025lipllm,
chaki2013probabilistic}.

\subsection{G2: Temporal Validity of Distributed State}
\label{sec:gap-time}

Distributed execution involves several distinct times: the time of the physical
event, the publication time, the arrival time, the time at which the consumer
can use the data, the target time of a forecast, and the end of the validity
interval.  Replacing the event time with a reception time stamp hides both
producer and network delay.  For a stream of state updates, the age of
information at time $t$ is
\begin{equation}
  \Delta(t)=t-u(t),
  \label{eq:aoi}
\end{equation}
where $u(t)$ is the event time of the newest update available to the consumer
\cite{kaul2012aoi}.  Age matters only relative to a decision: the same sample can
be fresh enough for a slow inspection task and too old for collision avoidance
between fast robots.  Comparing event times across robots also requires the
clocks to be comparable.  Synchronization protocols such as IEEE~1588 can bound
the offset between clocks \cite{ieee2020ptp}, but the bound helps the consumer
only if it travels with the time stamp.  Transport reliability or a cadence
notification does not establish that the margin a decision needs is available
\cite{omg2015dds,kronauer2021latency}.  A cross-robot interface therefore needs
event-time semantics, the clock domain and an error bound, an
application-specific validity bound, and a response to stale or incomparable
state.  \Cref{sec:temporal-admission} derives these fields from the collision
constraint in \eqref{eq:coupled-dynamics}.

\subsection{G3: Uncertainty, Provenance, and Evidence}

The word \emph{uncertainty} may refer to a covariance, an interval, a reachable
set, a calibrated probability, or semantic ambiguity.  These representations
cannot be compared or propagated in the same way, and a confidence value with no
declared meaning cannot satisfy a downstream chance constraint or robust margin.
The interface must identify the representation, the covered variables and
correlations, the calibration regime, the validity horizon, and the propagation
rule.

The same requirement applies to evidence.  A consumer should be able to tell a
proof from a solver status, a monitor verdict, a calibration report, a test
identifier, and a deployment measurement.  Provenance links a payload and its
evidence to the model, configuration, observations, and transformations that
produced them.  Without this lineage, a downstream rejection may reveal that a
claim failed but not which assumption or artifact needs repair
\cite{kim2026safeembodied,w3c2013prov}.

\subsection{G4: Implementation and Lifecycle Compatibility}

Type and wire compatibility are weaker than behavioral compatibility.  Two
components can decode the same message while using different coordinate frames,
units, parameter sets, middleware implementations, dependency versions, or clock
sources.  Admitting a result therefore requires separate identities for the
schema, protocol, implementation, dependencies, and configuration, together with
evidence that the combination is supported
\cite{macenski2022ros2,omg2015dds,kronauer2021latency}.

Message contracts also omit lifecycle unless it is represented explicitly.
Join, leave, restart, and partition recovery can invalidate old assignments or
duplicate resource ownership.  A multi-robot contract needs a team epoch,
membership and ownership state, legal lifecycle transitions, and a
reconciliation rule for retained or delayed data
\cite{luckcuck2026compositional,liu2026phyagentos}.

\subsection{G5: Intervention, Recovery, and Attribution}

Monitoring is incomplete if a failed check has no defined consequence.  The
responsible authority may request clarification, obtain fresh state,
re-estimate, replan, transfer control, degrade the mission, isolate a robot,
stop, or escalate to a human.  In a team, this response must be coordinated: an
unreported local safety override can invalidate peer trajectories or
task-completion logic.  Runtime-assurance and field-testing work accordingly
treats intervention, observability, and reporting as first-class mechanisms
\cite{desai2019soter,shivakumar2020soterros,caldas2024rvft}.

Fault origin and first observation must also be kept apart.  A collision risk
first detected by a controller may originate in semantic grounding, prediction,
physical modeling, communication delay, or hardware.  Recording origin, first
observation, effect, and recovery authority as separate fields supports both
online handling and post-mission diagnosis.

\subsection{Evidence from ROS~2 Middleware}
\label{sec:middleware-evidence}

Public failure reports from the nightly continuous-integration (CI) pipeline of
Nav2, the ROS~2 navigation stack, illustrate G2 and G4 in deployed software
\cite{nav2issue6163}.  Unchanged Nav2 test sources failed for different reasons
under three ROS middleware (RMW) implementations.
\begin{itemize}
  \item With \texttt{rmw\_zenoh}, the ROS~2 middleware based on Zenoh
    \cite{corsaro2023zenoh}, tests passed every assertion yet aborted at
    process exit: \texttt{rclcpp::shutdown()} ran from the destructor of a
    global test-fixture object after the thread-local storage of the Zenoh
    runtime had already been destroyed.  Moving initialization and shutdown
    into \texttt{main()} removed the failures \cite{nav2pr4948,nav2pr6241}.
  \item With Cyclone DDS, a teardown assertion in \texttt{ddsi\_fini} appeared
    with release~11.0.1 but not with 11.0.0, using the same
    \texttt{rmw\_cyclonedds\_cpp} version.  Reverting one upstream commit that
    changed the type-lookup request and reply headers eliminated the failure in
    40 repeated runs, although that commit may expose rather than contain the
    defect \cite{cyclonedds2426}.
  \item With Fast DDS, a service request reached its server and the callback
    ran, but the response was dropped because the response writer had not yet
    matched the client's response reader; increasing the client's wait time did
    not help.  The behavior is consistent with a known limitation of service
    discovery in the Fast DDS middleware layer
    \cite{nav2issue6163,fastrtps392}.
\end{itemize}

In each case, the message types, the application code, and the declared QoS
were unchanged.  What differed was the implementation identity, its patch
version, or the order of lifecycle events, which is exactly the information
that a type-level interface does not carry (G4).  The Fast DDS case also shows
that the availability of one endpoint is not evidence that a complete
request--response interaction can occur in time (G2).  The author contributed the
Zenoh fix and the Cyclone DDS diagnosis; all observations are taken from the
public record.  These are single-host test failures rather than multi-robot
field incidents, so they illustrate the gaps but do not measure their
prevalence.
